\chapter{Аналитическая часть}

В данной части будут рассмотрены последовательный и параллельный варианты алгоритма генерации изображения на основе трассировки лучей. Также будут представлены возможные подходы к распараллеливанию программы и осуществлён выбор одного из них.

\section{Трассировка лучей}

Трассировка лучей математически моделирует поведение света и его взаимодействие с физическими объектами --- отражение, преломление и тени. Сложные методы трассировки лучей позволяют создавать реалистичные спецэффекты для фильмов и игр~\cite{lit1}.

Из каждого пикселя камеры испускается луч. При попадании этого луча на какую либо поверхность, он отражается с учётом материала поверхности. 
Данный процесс может быть рекурсивным, что позволяет визуализировать
переотражения любой сложности.

Алгоритм трассировки лучей выполняется для каждого пикселя сцены. Из-за присутствия материалов с шершавой поверхностью возникает элемент случайности, в результате которого картинка становится шумной. Чтобы избавиться от этого недостатка необходимо запускать луч несколько раз, и усреднять полученные результаты. К тому же сам алгоритм является достаточно трудоёмким с точки зрения временной эффективности.

Однако, алгоритм подходит для распараллеливания, так как изображение можно разделить на отдельные участки, например на области или на сканирующие строки. Участки не зависят друг от друга, так как в расчёте не происходит изменения данных, которые читают разные потоки. Таким образом отпадает необходимость использовать методы синхронизации потоков.

\section{Параллельное исполнение программы}

Параллельное исполнение программы или её частей можно организовать несколькими способами:
\begin{itemize}
	\item зелёные потоки.
	\item нативные потоки;
	\item мультипроцессинг.
\end{itemize}

\subsection{Зелёные потоки}

Зелёные потоки ---  это легковесные потоки выполнения, создаваемые и управляемые в пользовательском пространстве средой выполнения (например, виртуальной машиной или библиотекой), а не ядром операционной системы. В отличие от нативных потоков, планирование и переключение между зелёными потоками осуществляется без прямого участия ОС.

Зелёные потоки эффективнее при создании и переключении, эта особенность делает их полезными в системах с высокой конкуренцией задач. Однако их ключевое ограничение заключается в том, что они не могут использовать несколько ядер процессора для истинного параллелизма: поскольку операционная система видит только один нативный поток, все зелёные потоки выполняются конкурентно, но не параллельно.  Кроме того, блокирующая операция в одном зелёном потоке может остановить выполнение всех остальных.

\subsection{Многопроцессная обработка}

Многопроцессная обработка --- это модель выполнения программ, при которой вычислительная нагрузка распределяется между несколькими независимыми процессами, каждый из которых обладает собственным адресным пространством и ресурсами. В отличие от многопоточности, где потоки разделяют память внутри одного процесса, процессы в многопроцессной обработке изолированы друг от друга, а их взаимодействие осуществляется с помощью механизмов межпроцессного взаимодействия, таких как: каналы, очереди сообщений, сокеты или разделяемая память. Эта модель эффективна на многоядерных процессорах, так как операционная система может распределять процессы по разным ядрам центрального процессора. Основные преимущества многопроцессной обработки:
\begin{itemize}
	\item возможность полного использования аппаратного параллелизма;
	\item высокая отказоустойчивость (сбой в одном процессе, не затрагивает другие);
	\item отсутствует необходимость синхронизировать доступ к памяти.
\end{itemize}

В то же время мультипроцессинг имеет и существенные недостатки: 
\begin{itemize}
	\item создание и переключение между процессами сопряжено со значительными накладными расходами по сравнению с потоками;
	\item обмен данными между процессами требует явного использования IPC-механизмов, что усложняет разработку и снижает производительность;
	\item каждый процесс потребляет дополнительную память, так как содержит собственную копию исполняемого кода и данных.
\end{itemize}

\subsection{Нативные потоки}
Поток --- это базовая единица выполнения в рамках процесса, представляющая собой последовательность инструкций, которая может выполняться независимо от других потоков того же процесса. Все потоки, принадлежащие одному процессу, разделяют его адресное пространство, включая код программы, глобальные переменные и открытые ресурсы, но при этом обладают собственными стеками и состояниями выполнения. Такая модель позволяет эффективно реализовывать конкурентность. На многоядерных системах потоки могут распределяться по разным ядрам, обеспечивая истинный параллелизм.

Основное преимущество потоков --- низкие накладные расходы по сравнению с процессами. Создание, переключение и взаимодействие между потоками происходит значительно быстрее, поскольку они используют общую память и не требуют сложных механизмов межпроцессного взаимодействия. Это делает многопоточность особенно эффективной для задач, требующих интенсивного обмена данными. Однако именно разделяемая память порождает ключевые проблемы многопоточного программирования: состояния гонки и взаимные блокировки, возникающие при некорректном доступе к общим данным. Для их предотвращения необходимо использовать примитивы синхронизации, такие как: мьютексы, семафоры, условные переменные. Это усложняет код и снижает его производительность.

\subsection{Выбор метода распараллеливания программы}

Для распараллеливания программы в лабораторной работе были выбраны нативные потоки, так как мультипроцессинг плохо подходит ввиду сложности организации доступа к разделяемой памяти, а использование зелёных потоков запрещено по условию работы.

 \section*{Вывод}
 
В данной части были представлены последовательный и параллельный варианты алгоритма трассировки лучей, а также выбран подход к организации параллельного выполнения программы.
\clearpage
